\section{Limitations}
\label{sec:conclusion:limitations}

The Phase~One work has six limitations that bear on the validity of
its conclusions and on the planning for Phase~Two. The first four
concern the dataset and the analytical claims it can support; the last
two concern the scope of the present work as a piece of design
research.

\begin{enumerate}
    \item \textbf{Calendar mismatch between the two data sources.}
    The AFC stream covers 1~January to 31~March 2026; the GPS stream
    covers 20~March to 10~May 2025. No calendar day is common to the
    two. Row-level joining between the two sources is therefore
    impossible, and the proposed system reconciles them only at the
    aggregate grain of route by hour by day type. Any claim made on
    the joined data is valid only at that grain.

    \item \textbf{Local GPS sample limitation.} The locally available
    GPS data covers a single day (18~March 2025) across nine routes;
    the full fifty-day window is documented in the dataset
    release~\cite{hcmcgps2025} but is not loaded locally for
    Phase~One. Several charts in the twenty-canvas catalogue depend
    on the full GPS window and carry a Phase-Two prerequisite status
    until the full dataset is loaded in P9 of the Phase~Two roadmap.

    \item \textbf{Absence of alighting events in the AFC stream.}
    The AFC stream records boarding events only. Full passenger
    journeys (origin to destination) cannot be reconstructed because
    no alighting event is recorded; live vehicle occupancy at a given
    moment is not directly computable. The proposed system surfaces a
    cumulative-boardings-per-trip proxy for crowding rather than a
    direct passenger count, with the proxy explicitly labelled.

    \item \textbf{Quality-of-Service indicator sensitivity to GPS
    noise.} The QoS-2 indicator (on-path operation, computed against
    a 50-metre route buffer) is systematically lower in dense urban
    areas because GPS positioning noise causes pings to fall outside
    the buffer even when the bus is on-route. The indicator is most
    informative when compared across routes rather than against an
    absolute threshold; the proposed system surfaces it as a relative
    metric and documents the caveat next to the indicator.

    \item \textbf{Generalisability beyond Ho Chi Minh City.} The design
    is bound to the specific shape of the HCMC bus dataset: the
    Vietnamese student-fare class HS-SV anchors the School-Aware
    Demand flagship, the Tết and 30~April calendar tagging is wired
    into the date dimension, and the operator and stop master tables
    follow the schema released by the city's transport authority.
    Portability to another Vietnamese city (Hanoi, {\DJ}\`a N\~\u{a}ng)
    would inherit the warehouse and visualisation rubric without
    change, but the Silver-tier dimension tables, the school
    geocoding ground truth, and the per-shock calendar windows would
    need to be rebuilt against the new city's data; the present work
    does not analyse the cost of that port.

    \item \textbf{Self-audit rather than external validation.} The
    Phase-One requirements evaluation of
    Section~\ref{sec:impl:eval} is a self-audit against the
    specification this report itself authors. External validation in
    the form of advisor sign-off, stakeholder review of the fourteen
    stakeholder questions, and the two-wave user evaluation across
    the three user roles is scheduled for Phase~Two
    (Section~\ref{sec:future:eval}). No claim of external validation
    is made in Phase~One.
\end{enumerate}
